\chapter{总结与展望}

\section{工作总结}
本文围绕《On the Optimal Time/Space Tradeoff for Hash Tables》提出的结构思想，完成了其在工程环境中的复现与系统化实现。本文实现通过 bin/mini-bin/fallback 的分层布局，使得查询仅需访问偏好 mini-bin 与 fallback 两段连续下标区间，从而在“区间扫描次数”意义下将查询探针限制为常数。插入方面，实现了有界深度的 k-kick 置换策略与 fallback 兜底路径，并在数据库化场景下通过“锁定单个 bin + 内存模拟踢出链 + 批量写回”的方式减少 SQL 往返。系统层面，本文提供了完整的 HTTP 接口与前端实验面板，并实现在线扩容流程以支撑高负载下的长期运行与可视化分析。

\section{不足与改进方向}
数据库化实现在可复现与可观测方面具有优势，但也带来额外开销与工程约束。后续可从以下方向改进：
\begin{itemize}
  \item \textbf{性能优化}：进一步减少范围查询与事务开销，例如更细粒度的批量读写、更激进的缓存策略或预取；
  \item \textbf{并发与隔离}：在扩容期间的读写隔离、跨线程与跨进程的扩容状态持久化，以及更稳健的失败恢复；
  \item \textbf{实验完善}：补齐不同分布（uniform/偏斜）与不同参数（\(k\)、load factor、mini-bin 派生规则）下的全面对比，并与经典结构进行对照实验；
  \item \textbf{功能扩展}：支持更多键值类型、批量删除与重建策略，以及更丰富的可视化与数据导出能力。
\end{itemize}

